
\begin{figure}[H]
    \caption{Delimitación contextual del estudio}
    \label{fig:delimitacion-contextual}

    \centering

    \begin{tikzpicture}[
        bloque/.style={
            draw=black!60,
            rounded corners=2pt,
            align=center,
            text width=5cm,
            minimum height=1.15cm,
            inner sep=8pt,
            font=\small
        },
        flecha/.style={
            -{Stealth[length=2.5mm]},
            thick,
            black!65
        }
    ]

        % Delimitación geográfica y poblacional
        \node[bloque] (ciudad) {
            \textbf{Bogotá, D. C.}
        };

        \node[bloque, below=0.35cm of ciudad] (instituciones) {
            Instituciones de educación superior
        };

        \node[bloque, below=0.35cm of instituciones] (poblacion) {
            Población propuesta: estudiantes de 18 años o más
        };

        % Objeto de estudio
        \node[
            bloque,
            text width=4.2cm,
            minimum height=2.2cm,
            right=1.3cm of instituciones
        ] (objeto) {
            \textbf{Objeto de estudio}\\[4pt]
            Herramientas digitales para fines académicos
        };

        % Relaciones conceptuales
        \draw[flecha] (ciudad) -- (instituciones);
        \draw[flecha] (instituciones) -- (poblacion);
        \draw[flecha] (instituciones.east) -- (objeto.west);

    \end{tikzpicture}

    \par\smallskip
    \begin{minipage}{0.9\textwidth}
        \footnotesize
        \textit{Nota.} Esquema conceptual de delimitación, no mapa cartográfico ni representación de la ubicación exacta de instituciones.
    \end{minipage}
\end{figure}
